\section{Threats to Validity}

\subsection{Internal Validity}
A portion of the validation uses scripted southbound interactions and emulation-assisted workflows. This is appropriate for verifying protocol handling, orchestration semantics, and status convergence, but it may hide some low-level hardware effects that would emerge on physical boards under noisy environmental conditions or tighter timing constraints. To mitigate this threat, the experiments focus their claims on system-level functional correctness rather than on precise real-time performance.

\subsection{External Validity}
The presented deployment uses a single-node K3s cluster and a bounded experimental scenario. Larger deployments with multiple gateways, geographically distributed links, or mixed wireless conditions may introduce failure modes and scheduling behaviors not exercised here. The results should therefore be interpreted as evidence that the design is viable and reproducible, not as a complete characterization of all operational environments.

\subsection{Construct Validity}
The study emphasizes end-to-end correctness, secure attachment, liveness observability, and transaction-oriented workflow completion. These constructs align with the stated contribution, but they do not cover every dimension a production platform would need. In particular, energy consumption, memory footprint on physical boards, multi-device contention, and large-scale fault tolerance are outside the present evaluation scope. Moreover, some quantitative decisions rely on authoritative runtime signals from the gateway and API server rather than on CRD phase alone, because the current artifact still exposes noisy replicated state in part of the northbound status path. Future work should close that gap so that the control plane itself becomes the unambiguous source of truth.

\subsection{Conclusion Validity}
Because the paper reports a full implementation and not only a simulation, the conclusions are grounded in actual deployment behavior. Nevertheless, the argument remains strongest for the claim that gateway-centric orchestration can be made operationally coherent on constrained targets. Claims about optimality, universal superiority, or production readiness at industrial scale would require broader comparative evaluation and long-duration testing, which are intentionally left outside the current manuscript.
